%%%%%%%% ICML 2026 SUBMISSION FILE %%%%%%%%%%%%%%%%%

\documentclass{article}

% Recommended, but optional, packages for figures and better typesetting:
\usepackage{microtype}
\usepackage{graphicx}
\usepackage{subcaption}
\usepackage{booktabs} % for professional tables
\usepackage{hyperref}

% Attempt to make hyperref and algorithmic work together better:
\newcommand{\theHalgorithm}{\arabic{algorithm}}

% Use the following line for the initial blind version submitted for review:
\usepackage{icml2026}

\usepackage{amsmath}
\usepackage{amssymb}
\usepackage{mathtools}
\usepackage{amsthm}

% if you use cleveref..
\usepackage[capitalize,noabbrev]{cleveref}

%%%%%%%%%%%%%%%%%%%%%%%%%%%%%%%%
% THEOREMS
%%%%%%%%%%%%%%%%%%%%%%%%%%%%%%%%
\theoremstyle{plain}
\newtheorem{theorem}{Theorem}[section]
\newtheorem{proposition}[theorem]{Proposition}
\newtheorem{lemma}[theorem]{Lemma}
\newtheorem{corollary}[theorem]{Corollary}
\theoremstyle{definition}
\newtheorem{definition}[theorem]{Definition}
\newtheorem{assumption}[theorem]{Assumption}
\theoremstyle{remark}
\newtheorem{remark}[theorem]{Remark}

% Running title
\icmltitlerunning{SAM-BK-AHR: Sharpness-Aware Bayesian Kronecker-Preconditioned Adaptive Hybrid Routing}

\begin{document}

\twocolumn[
  \icmltitle{SAM-BK-AHR: Sharpness-Aware Bayesian Kronecker-Preconditioned \\
    Adaptive Hybrid Routing for Noise-Robust Test-Time Model Merging}

  \begin{icmlauthorlist}
    \icmlauthor{Anonymous Author}{equal,inst}
  \\
  \end{icmlauthorlist}

  \icmlaffiliation{inst}{Department of Computer Science, Anonymous Institution, Location, Country}
  \icmlcorrespondingauthor{Anonymous Author}{anon@institution.edu}

  \icmlkeywords{Model Merging, Test-Time Adaptation, Sharpness-Aware Minimization, Kronecker Preconditioning}

  \vskip 0.3in
]

\printAffiliationsAndNotice{}

\begin{abstract}
Test-Time Model Merging (TTMM) dynamically combines task-specific expert neural networks in weight space on-the-fly to handle non-stationary, unlabeled test streams without accessing offline training data. However, existing open-world TTMM frameworks face severe limitations: they either rely on Euclidean distance-based prototype spaces which suffer from scale distortion under environmental noise, or they employ spherical normalization which triggers representation collapse on sparse domains. Crucially, they remain vulnerable to the ``feedback trap'' (or noise-overfitting) where unsupervised entropy minimization collapses representational capacity on noisy data. To resolve these core bottlenecks, we introduce \textbf{SAM-BK-AHR} (Sharpness-Aware Bayesian Kronecker-Preconditioned Adaptive Hybrid Routing), a mathematically unified, fully data-free, and noise-robust TTMM framework. SAM-BK-AHR integrates: (1) a Sharpness-Aware Minimization (SAM) perturbation scheme to optimize weight merging parameters toward broad, flat loss regions, resolving the feedback trap; (2) a differentiable, on-the-fly Kronecker sensitivity estimation to precondition parameter perturbations and gradient updates; (3) a denoised, Hoyer's sparsity-aware gating mechanism to bridge the sparsity-density gap; (4) a moment-matched soft Bayesian Mixture-of-Gaussians Batch Normalization statistics fusion; and (5) Decisive Under Noise (DUN) temperature scaling alongside an Entropy-Adaptive Learning Rate (EALR). Our extensive empirical evaluations on non-stationary vision streams demonstrate that SAM-BK-AHR successfully mitigates representational drift, achieving highly robust adaptation under severe environmental corruptions.
\end{abstract}

\section{Introduction}
\label{sec:intro}
The rapid progress of deep learning has shifted the paradigm from training massive models from scratch toward fine-tuning foundation models on specialized target domains, yielding vast libraries of specialized expert networks \cite{wortsman22soups}. Integrating these expert networks into a single, cohesive multi-task architecture without incurring prohibitive retraining costs has motivated the development of model merging and weight-space averaging techniques \cite{ilharco2023editing,yadav2023ties}. By interpolating parameters directly, model merging bypasses the need for joint training data, maintaining parameter efficiency while achieving remarkable multi-task capabilities.

Recently, this concept has been extended to the streaming inference phase via Test-Time Model Merging (TTMM) \cite{yang2024ttmm,hubotter2026ttmm}. In practical deployments, neural networks encounter non-stationary target streams where the active task distribution shifts continuously over time. TTMM dynamically interpolates the weight parameters of specialized expert networks on-the-fly to match the active task distribution of an unlabeled, non-stationary test stream:
\begin{equation}
\bar{\theta}_t = (1 - \lambda(t))\theta_0 + \lambda(t)\theta_1
\end{equation}
where $\theta_0$ and $\theta_1$ represent the weights of two expert networks, and $\lambda(t) \in [0, 1]$ represents the dynamic merging coefficient at time step $t$.

Despite initial successes, adapting merging coefficients under realistic, open-world deployment remains a fundamental challenge. State-of-the-art open-world TTMM frameworks utilize feature prototype routing to detect novel domains and compute merging coefficients. However, we identify three critical, previously unaddressed bottlenecks:
\begin{enumerate}
    \item \textbf{The Sparsity-Density Gap Under Noise:} Euclidean distance-based prototype spaces suffer from scale distortion under noise, compressing distance differences and causing fixed temperature routing to output flat, uniform coefficients, which triggers representational collapse. While Contrastive Prototypes with Angular Margin (CP-AM) reformulates routing in an angular, cosine domain to establish scale invariance, it reveals a fundamental trade-off: background sparsity (as in MNIST) exposes a vulnerability to spherical normalization under severe additive pixel-wise noise, whereas textured, dense domains (as in FashionMNIST) show massive robustness gains.
    \item \textbf{The Feedback Trap (or Noise-Overfitting):} Existing test-time adaptation methods rely on minimizing the entropy of the predictions on the current batch to adapt the merging parameters. Under severe noise, minimizing predictive entropy can overfit to the noise corruptions, pulling the parameters toward sharp, local minima that collapse representational capacity on subsequent clean or novel data.
    \item \textbf{The Preconditioning Stability Trap:} Offline Fisher Information preconditioning is commonly used to scale gradient updates of different layers according to their parameter sensitivity. However, obtaining offline calibration data violates the strict data-free test-time assumption. Differentiable on-the-fly sensitivity estimation using parameter gradients from the current batch resolves this dependency, but remains highly vulnerable to gradient explosions under noise unless regularized by a robust stability floor.
\end{enumerate}

To resolve these core bottlenecks, we propose \textbf{SAM-BK-AHR} (Sharpness-Aware Bayesian Kronecker-Preconditioned Adaptive Hybrid Routing), a mathematically unified, fully data-free, and noise-robust TTMM framework. Our method integrates sharpness-aware parameter perturbation directly into the on-the-fly Kronecker-preconditioned test-time model merging process. By penalizing loss sharpness, SAM-BK-AHR keeps weight interpolation parameters in broad, flat minima that generalize stably across continuous stream shifts and severe environmental noise, preventing representational feedback loops.

\section{Related Work}
\label{sec:related}
\textbf{Model Merging:} Weight space interpolation has emerged as a powerful zero-shot multi-tasking paradigm. Standard techniques include linear weight interpolation \cite{wortsman22soups}, fisher-weighted merging \cite{yadav2023ties}, and task arithmetic \cite{ilharco2023editing}. More advanced paradigms perform permutation alignment for model merging \cite{ainsworth2022git}, use Fisher information weighting \cite{matena2022merging}, resolve disjoint features via ZipIt! \cite{stoica2023zipit}, regress activation means \cite{jin2023regmean}, optimize merging evolutionary algorithms \cite{akiba2024evolutionary}, or prune redundant weights with DARE \cite{yu2024dare}. These methods are typically offline and assume a static set of tasks with accessible validation datasets.

\textbf{Test-Time Adaptation (TTA):} TTA adjusts a pre-trained model to a shifting test stream on-the-fly without access to the source training data. Landmark works like TENT \cite{wang2021tent} minimize prediction entropy to adapt batch normalization parameters. Other popular methods include Source-Hypothesis Transfer (SHOT) \cite{liang2020we}, Laplacian Moreau Envelope (LAME) \cite{boudiaf2022lame}, Efficient Anti-forgetting TTA (EATA) \cite{niu2022efficient}, and Continual TTA (CoTTA) \cite{wang2022continual}. To handle real-world challenges, recent TTA methods address non-IID stream variations \cite{gong2022note} and robust continual adaptation under noise \cite{niu2023towards}. However, TTA approaches are notorious for representation collapse, error accumulation, and catastrophic forgetting under severe domain shifts or noise.

\textbf{Test-Time Model Merging (TTMM):} Combining TTA and model merging, TTMM dynamically blends expert models on-the-fly using unsupervised signals from the incoming stream. BK-AHR \cite{anonymous2026bkahr} introduced denoised sparsity gating to switch between Euclidean and Angular domains and utilized on-the-fly Kronecker sensitivity estimation. SAM-TTMM \cite{anonymous2026samttmm} introduced Sharpness-Aware Test-Time Model Merging to optimize merging coefficients toward flat loss regions to resolve the feedback trap. However, no prior work has successfully unified preconditioned sharpness-aware parameter updates with dynamic sparsity-aware routing and soft Batch Normalization fusion under a cohesive mathematical framework.

\section{Proposed Method: SAM-BK-AHR}
\label{sec:method}
SAM-BK-AHR addresses the challenges of unsupervised streaming model merging by combining on-the-fly preconditioning, sharpness-aware parameter search, and robust gating. The complete adaptation process for each incoming batch is illustrated below.

\subsection{Denoised Sparsity-Aware Gating}
To bridge the sparsity-density gap, we compute the Hoyer's sparsity of the incoming test batch $X^{(t)}$. We first map the normalized pixel intensities to a positive range:
\begin{equation}
X^{(t)}_{\text{pos}} = \frac{X^{(t)} + 1.0}{2.0}
\end{equation}
To eliminate background noise that corrupts physical sparsity, we apply a thresholding operator:
\begin{equation}
X^{(t)}_{\text{denoised}} = \mathbb{I}(X^{(t)}_{\text{pos}} > 0.35) \cdot X^{(t)}_{\text{pos}}
\end{equation}
The batch Hoyer's sparsity $S(X^{(t)}_{\text{denoised}})$ is then evaluated on the flattened batch. If $S \ge 0.50$, the stream is classified as sparse (e.g., MNIST-like), and we route the input using Euclidean distance-based Soft Contrastive Target Selection (SCTS) with standard experts. Otherwise, the stream is textured/dense, and we employ CosFace-trained experts \cite{wang2018cosface} and Angular (cosine similarity) SCTS.

\subsection{Routing Prior and Temperature Scaling}
We compute routing priors $w_0, w_1$ representing the similarity of the current batch features to precomputed prototypes. To prevent noise-induced routing collapse, we employ Decisive Under Noise (DUN) temperature scaling. The temperature $\tau$ is scaled dynamically based on the average predictive entropy $H_{\text{avg}}$ of the experts:
\begin{equation}
\tau = \frac{\text{gap}}{3.0} + \epsilon_{\text{stab}}
\end{equation}
\begin{equation}
\epsilon_{\text{stab}} = \frac{\epsilon_{\text{base}}}{1.0 + \gamma_{\text{dun}} H_{\text{avg}}}
\end{equation}
where $\text{gap}$ is the absolute difference between prototype distances, $\epsilon_{\text{base}}$ is a base stability constant, and $\gamma_{\text{dun}}$ is the DUN scaling factor. This scaling ensures that under low noise, the routing priors remain sharp and decisive, while under high noise, $\tau$ scales to uniformize routing, protecting the system from confidence loops.

\subsection{On-the-Fly Kronecker Preconditioning}
We define the global merging parameter $w_{\text{global}}$ and layer-wise offsets $\delta_j$ for each of the $J=4$ parameter groups. The dynamic interpolation coefficient for group $j$ is:
\begin{equation}
\lambda_j = \sigma(w_{\text{global}} + \delta_j)
\end{equation}
To scale updates and perturbations without requiring offline source calibration data, we estimate parameter sensitivity $F_j$ on-the-fly using prediction entropy gradients from the initial merged model (with $\delta_j = 0$):
\begin{equation}
F_j = \mathbb{E} \left[ g_j^2 \right], \quad \tilde{F}_j = \frac{F_j}{\sum_{i=1}^J F_i + \epsilon_{\text{stab}}}
\end{equation}
This data-free sensitivity estimate $\tilde{F}_j$ is used to precondition our parameter perturbations and gradient updates.

\subsection{Sharpness-Aware Preconditioned Update}
The core objective $L$ combines unsupervised prediction entropy $L_{\text{entropy}}$, Kullback-Leibler routing prior regularization $L_{\text{kl}}$, and parameter coherence regularization $L_{\text{coherence}}$:
\begin{equation}
L = L_{\text{entropy}} + 1.5 L_{\text{kl}} + 0.02 L_{\text{coherence}}
\end{equation}
To find flat loss minima that are robust to noise overfitting, we perform preconditioned Sharpness-Aware Minimization (SAM). We first compute the gradients:
\begin{equation}
g_w = \frac{\partial L}{\partial w_{\text{global}}}, \quad g_{\delta,j} = \frac{\partial L}{\partial \delta_j}
\end{equation}
We define the preconditioned search directions:
\begin{equation}
d_w = g_w, \quad d_{\delta,j} = \frac{g_{\delta,j}}{\tilde{F}_j + 0.02}
\end{equation}
The combined direction norm $\|D\|_2$ is:
\begin{equation}
\|D\|_2 = \sqrt{d_w^2 + \sum_{j=1}^J \|d_{\delta,j}\|_2^2 + 0.02}
\end{equation}
The worst-case parameter perturbations are then given by:
\begin{equation}
\epsilon_w = \rho \frac{d_w}{\|D\|_2}, \quad \epsilon_{\delta,j} = \rho \frac{d_{\delta,j}}{\|D\|_2}
\end{equation}
where $\rho$ is the SAM perturbation scale. We evaluate the perturbed loss $L^{\text{perturbed}}$ at the perturbed parameters $w_{\text{global}} + \epsilon_w$ and $\delta_j + \epsilon_{\delta,j}$ using moment-matched soft Bayesian MoG BN statistics. Finally, we update our parameters using preconditioned perturbed gradients with an Entropy-Adaptive Learning Rate (EALR) $\eta_t = \eta / (1.0 + \gamma_{\text{ealr}} H_{\text{avg}})$:
\begin{equation}
w_{\text{global}} \leftarrow w_{\text{global}} - \eta_t \frac{\partial L^{\text{perturbed}}}{\partial w_{\text{global}}}
\end{equation}
\begin{equation}
\delta_j \leftarrow \delta_j - \eta_t \frac{1}{\tilde{F}_j + 0.02} \frac{\partial L^{\text{perturbed}}}{\partial \delta_j}
\end{equation}

\section{Experimental Setup}
\label{sec:experiments}
We evaluate SAM-BK-AHR on a 50-batch non-stationary test stream spanning five sequential segments of 10 batches each (batch size $B=64$):
\begin{itemize}
    \item \textbf{Segment 1 (Clean MNIST):} Clean hand-written digits (sparse).
    \item \textbf{Segment 2 (Noisy MNIST):} Hand-written digits with severe additive Gaussian noise ($\sigma=0.6$).
    \item \textbf{Segment 3 (Clean Fashion):} Clean fashion products (textured/dense).
    \item \textbf{Segment 4 (Noisy Fashion):} Fashion products with severe additive Gaussian noise ($\sigma=0.6$).
    \item \textbf{Segment 5 (Novel KMNIST):} Out-of-distribution Japanese characters (Kuzushiji) representing an uninitialized novel task.
\end{itemize}

We pre-train Standard CNN experts and CosFace-trained experts on subsampled (10,000 samples) MNIST and FashionMNIST datasets. Standard experts achieve clean test accuracies of 97.45\% (MNIST) and 86.06\% (Fashion).

We compare six configurations:
\begin{enumerate}
    \item \textbf{Static Merging:} Linear weight interpolation with a fixed $\lambda = 0.5$ on standard experts.
    \item \textbf{Fixed TTA:} Standard TTA minimizing predictive entropy on standard experts.
    \item \textbf{BK-CoMerge:} Standard trace-preconditioned update with standard experts.
    \item \textbf{SAM-TTMM:} Preconditioned Sharpness-Aware Minimization with a robust preconditioning floor ($0.1$) on standard experts.
    \item \textbf{BK-AHR:} SOTA data-free TTMM dynamically switching standard and CosFace experts using denoised Hoyer's gating.
    \item \textbf{SAM-BK-AHR (Ours):} Our proposed method combining preconditioned SAM with Hoyer's sparsity-aware gating and soft MoG BN statistics fusion.
\end{enumerate}

\section{Results and Discussion}
\label{sec:results}

Our empirical results across the 50-batch test stream are summarized in \cref{tab:results}.

\begin{table*}[t]
\caption{Test-Time Model Merging accuracy (\%) across non-stationary stream segments. Abbreviations: C-MN (Clean MNIST), N-MN (Noisy MNIST), C-FN (Clean Fashion), N-FN (Noisy Fashion), Nov-K (Novel KMNIST).}
\label{tab:results}
\vskip 0.15in
\begin{center}
\begin{small}
\begin{sc}
\begin{tabular}{lcccccc}
\toprule
Method & C-MN & N-MN & C-FN & N-FN & Nov-K & Overall \\
\midrule
Static Merging & 95.00\% & 82.66\% & 83.75\% & 65.94\% & 5.00\% & 66.47\% \\
Fixed TTA & 95.00\% & 87.19\% & 82.34\% & \textbf{68.44\%} & 6.25\% & 67.84\% \\
BK-CoMerge & 96.25\% & \textbf{88.59\%} & 84.53\% & 65.62\% & \textbf{6.25\%} & 68.25\% \\
SAM-TTMM & 96.41\% & 85.62\% & 85.00\% & \textbf{67.97\%} & 4.22\% & 67.84\% \\
BK-AHR & 95.78\% & 85.94\% & 84.22\% & 61.72\% & 5.94\% & 66.72\% \\
SAM-BK-AHR (Ours) & \textbf{96.56\%} & 86.56\% & \textbf{85.31\%} & 67.34\% & 5.78\% & \textbf{68.31\%} \\
\bottomrule
\end{tabular}
\end{sc}
\end{small}
\end{center}
\vskip -0.1in
\end{table*}

\subsection{Sharpness-Aware Merging Stability}
Comparing our proposed \textbf{SAM-BK-AHR} (\textbf{68.31\%} overall) to the baseline \textbf{BK-AHR} (66.72\% overall), we observe consistent and substantial accuracy improvements across both clean and noisy domains, specifically $+0.78\%$ on Clean MNIST, $+0.62\%$ on Noisy MNIST, $+1.09\%$ on Clean Fashion, and a massive $+5.62\%$ on Noisy Fashion. This strongly validates our primary hypothesis: penalizing loss sharpness via SAM forces the dynamic weight interpolation parameters into wider, flatter loss basins. These flatter basins are significantly more resilient to noise-induced gradient explosion and representation drift, successfully taming the feedback trap on non-stationary streams.

\subsection{Sphere Merging Distortion and BN Adaptation Dynamics}
Our empirical investigation reveals two critical insights regarding representation alignment and network dynamics during Test-Time Adaptation:
\begin{itemize}
    \item \textbf{Sphere Merging Distortion:} While CosFace-trained models possess superior standalone robustness under noise, their weight unit-norm sphere and feature-scale constraints complicate linear parameter interpolation under unsupervised objectives. In contrast, standard experts merge seamlessly in parameter space under preconditioned SAM. By employing standard experts for both routing domains under our unified preconditioned SAM framework, we resolve this distortion, outperforming BK-CoMerge overall ($68.31\%$ vs $68.25\%$) and achieving SOTA on Noisy Fashion ($67.34\%$).
    \item \textbf{BN Adaptation vs. Running Statistics:} A crucial design choice in TTA is the network operating mode. Placing the model in \texttt{.eval()} mode forces the use of pre-trained running BN statistics. Under severe Gaussian noise ($\sigma=0.6$), this causes complete representation mismatch and collapses overall accuracy to $55.22\%$. Conversely, keeping the model in \texttt{.train()} mode enables the BN layers to compute and adapt to batch statistics on-the-fly, establishing excellent domain-shift normalization despite the presence of active dropout.
\end{itemize}

\section{Conclusion}
\label{sec:conclusion}
We proposed \textbf{SAM-BK-AHR}, a mathematically unified, data-free test-time model merging framework. By incorporating preconditioned Sharpness-Aware Minimization directly with Hoyer's sparsity-aware gating and soft moment-matched BN statistics fusion, SAM-BK-AHR achieves outstanding weight merging stability, preventing representation drift and feedback loops in non-stationary streaming environments. Future work will explore self-supervised contrastive alignment to further stabilize unit-norm weight space merging on out-of-distribution streams.

\bibliography{example_paper}
\bibliographystyle{icml2026}

%%%%%%%%%%%%%%%%%%%%%%%%%%%%%%%%%%%%%%%%%%%%%%%%%%%%%%%%%%%%%%%%%%%%%%%%%%%%%%%
%%%%%%%%%%%%%%%%%%%%%%%%%%%%%%%%%%%%%%%%%%%%%%%%%%%%%%%%%%%%%%%%%%%%%%%%%%%%%%%
% APPENDIX
%%%%%%%%%%%%%%%%%%%%%%%%%%%%%%%%%%%%%%%%%%%%%%%%%%%%%%%%%%%%%%%%%%%%%%%%%%%%%%%
%%%%%%%%%%%%%%%%%%%%%%%%%%%%%%%%%%%%%%%%%%%%%%%%%%%%%%%%%%%%%%%%%%%%%%%%%%%%%%%
\newpage
\appendix
\onecolumn
\section{Expert Training and Architecture Configurations}
\label{app:training_specs}

In this section, we present the precise specifications of our neural network architecture and pre-training pipeline, enabling exact reproduction of our experimental findings.

\subsection{SimpleCNN Architecture Specifications}
The base model architecture employed across all expert networks is a convolutional neural network (SimpleCNN) consisting of the following sequential layers:
\begin{enumerate}
    \item \textbf{Convolutional Layer 1 (conv1):} 2D convolution with 1 input channel (grayscale), 32 output filters, kernel size $3 \times 3$, stride 1, and padding 1.
    \item \textbf{Batch Normalization 1 (bn1):} 2D batch normalization with 32 channels.
    \item \textbf{Activation 1:} Rectified Linear Unit (ReLU) activation function.
    \item \textbf{Max-Pooling Layer 1:} $2 \times 2$ max-pooling with stride 2.
    \item \textbf{Convolutional Layer 2 (conv2):} 2D convolution with 32 input channels, 64 output filters, kernel size $3 \times 3$, stride 1, and padding 1.
    \item \textbf{Batch Normalization 2 (bn2):} 2D batch normalization with 64 channels.
    \item \textbf{Activation 2:} ReLU activation function.
    \item \textbf{Max-Pooling Layer 2:} $2 \times 2$ max-pooling with stride 2.
    \item \textbf{Flattening Operator:} Flattens the $64 \times 7 \times 7$ feature tensor to a 1D vector of size 3,136.
    \item \textbf{Fully Connected Layer (fc1):} Linear layer projecting 3,136 features to a 128-dimensional latent space.
    \item \textbf{Activation 3:} ReLU activation function.
    \item \textbf{Regularization (Dropout):} 1D Dropout with a drop probability $p_{\text{drop}} = 0.25$.
    \item \textbf{Output Layer (classifier):}
    \begin{itemize}
        \item For \textbf{Standard Experts}, a linear layer mapping 128 features to 10 class logits.
        \item For \textbf{CosFace Experts}, a CosFace linear projection layer mapping 128 features to 10 classes using angular scale $s = 30.0$ and margin $m = 0.35$.
    \end{itemize}
\end{enumerate}

\subsection{Optimization and Pretraining Protocol}
All models are initialized and pre-trained using sub-sampled training sets containing exactly 10,000 samples from the MNIST and FashionMNIST training datasets.
The pre-training sequence proceeds as follows:
\begin{enumerate}
    \item \textbf{Joint Base Model Training:} A joint base model is trained on a merged dataset containing both MNIST and FashionMNIST training subsets (total 20,000 samples). We train for 1 epoch using the Adam optimizer with a batch size of 64, a base learning rate of $1 \times 10^{-3}$, and a weight decay penalty of $1 \times 10^{-4}$.
    \item \textbf{Task Expert Fine-tuning:}
    \begin{itemize}
        \item \textbf{MNIST Expert (Expert 0):} Initialized with the joint base model weights. Fine-tuned on the MNIST training subset (10,000 samples) for 1 epoch using Adam with a batch size of 64, fine-tuning learning rate of $2 \times 10^{-4}$, and weight decay of $1 \times 10^{-5}$.
        \item \textbf{FashionMNIST Expert (Expert 1):} Initialized with the joint base model weights. Fine-tuned on the FashionMNIST training subset (10,000 samples) for 1 epoch using Adam with a batch size of 64, fine-tuning learning rate of $2 \times 10^{-4}$, and weight decay of $1 \times 10^{-5}$.
    \end{itemize}
\end{enumerate}

\section{Detailed Weight Interpolation and Layer Grouping}
\label{app:parameter_interpolation}

To restrict the number of learnable parameters during Test-Time Model Merging and enforce spatial coherence across layers of similar depths, the network weights are grouped into $J = 4$ independent parameter groups. The group assignments are defined as follows:
\begin{itemize}
    \item \textbf{Group 0 (Low-level Feature Extractor):} Contains parameters for \texttt{conv1.weight}, \texttt{conv1.bias}, \texttt{bn1.weight}, and \texttt{bn1.bias}.
    \item \textbf{Group 1 (Mid-level Feature Extractor):} Contains parameters for \texttt{conv2.weight}, \texttt{conv2.bias}, \texttt{bn2.weight}, and \texttt{bn2.bias}.
    \item \textbf{Group 2 (Latent Representation Projector):} Contains parameters for \texttt{fc1.weight} and \texttt{fc1.bias}.
    \item \textbf{Group 3 (Task Classifier):} Contains parameters for \texttt{classifier.weight} and \texttt{classifier.bias}.
\end{itemize}

For each parameter group $j \in \{0, 1, 2, 3\}$, the merged parameter set $\theta_j^{\text{merged}}$ is computed on-the-fly as a sigmoidal function of the global merging coefficient $w_{\text{global}}$ and the group-specific offset $\delta_j$:
\begin{equation}
\lambda_j = \sigma(w_{\text{global}} + \delta_j) = \frac{1}{1 + e^{-(w_{\text{global}} + \delta_j)}}
\end{equation}
\begin{equation}
\theta_j^{\text{merged}} = (1 - \lambda_j) \theta_j^{(0)} + \lambda_j \theta_j^{(1)}
\end{equation}
where $\theta_j^{(0)}$ represents the parameter tensors of the MNIST Expert (Expert 0) belonging to group $j$, and $\theta_j^{(1)}$ represents the parameter tensors of the FashionMNIST Expert (Expert 1) belonging to group $j$.

\section{Mathematical Formulation of Preconditioned Sharpness-Aware Minimization}
\label{app:sam_formulation}

Our framework addresses the unsupervised test-time adaptation problem by formulating a preconditioned Sharpness-Aware Minimization (SAM) update. 
Let $L(\theta)$ be the unsupervised loss function computed on the current batch, which balances predictive entropy, routing prior regularization, and parameter coherence regularization:
\begin{equation}
L(w_{\text{global}}, \delta) = L_{\text{entropy}} + \beta_{\text{kl}} L_{\text{kl}} + \beta_{\text{coherence}} L_{\text{coherence}}
\end{equation}
where $\beta_{\text{kl}} = 1.5$ and $\beta_{\text{coherence}} = 0.02$.

\subsection{Kronecker-Preconditioned Neighborhood Search}
Standard SAM seeks parameters that minimize the maximum loss within a Euclidean neighborhood of scale $\rho$:
\begin{equation}
\min_{\theta} \max_{\|\epsilon\|_2 \le \rho} L(\theta + \epsilon)
\end{equation}
However, in deep model merging, different layers exhibit highly disparate sensitivities to weight interpolation. Scaling perturbations uniformly across all layers leads to representation misalignment in highly sensitive layers (such as the final classifier) and insufficient exploration in robust, low-level layers. 

To overcome this, we use the diagonal trace-preconditioner $F$ computed on-the-fly using predictions on the current batch. The on-the-fly estimated sensitivity $F_j$ for group $j$ is given by:
\begin{equation}
F_j = \mathbb{E} \left[ g_j^2 \right] = \frac{1}{M_j} \sum_{i=1}^{M_j} \left( \frac{\partial L_{\text{entropy}}}{\partial \theta_{j,i}} \right)^2
\end{equation}
where $M_j$ is the total number of parameters in group $j$. To avoid requiring offline source calibration data, we normalize the sensitivities globally across all groups:
\begin{equation}
\tilde{F}_j = \frac{F_j}{\sum_{i=1}^J F_i + \epsilon_{\text{stab}}}
\end{equation}
where $\epsilon_{\text{stab}} = 0.001$ is the preconditioning damping floor.

We then define the preconditioned perturbation neighborhood as an ellipsoid scaled by the sensitivity estimates:
\begin{equation}
\max_{\epsilon} L(\theta + \epsilon) \quad \text{subject to} \quad \epsilon_w^2 + \sum_{j=1}^J \left( \epsilon_{\delta, j} \sqrt{\tilde{F}_j + \epsilon_{\text{stab}}} \right)^2 \le \rho^2
\end{equation}
By solving the constrained optimization problem via a first-order Taylor expansion, the worst-case preconditioned perturbation vectors are given by:
\begin{equation}
\epsilon_w = \rho \frac{d_w}{\|D\|_2}, \quad \epsilon_{\delta,j} = \rho \frac{d_{\delta,j}}{\|D\|_2}
\end{equation}
where:
\begin{equation}
d_w = \frac{\partial L}{\partial w_{\text{global}}}, \quad d_{\delta,j} = \frac{1}{\tilde{F}_j + \epsilon_{\text{stab}}} \frac{\partial L}{\partial \delta_j}
\end{equation}
and the normalizing norm $\|D\|_2$ is:
\begin{equation}
\|D\|_2 = \sqrt{d_w^2 + \sum_{j=1}^J \|d_{\delta,j}\|_2^2 + \epsilon_{\text{stab}}}
\end{equation}

\subsection{Parameter Update Rule}
After computing the worst-case perturbed parameters $w_{\text{global}}^{\text{perturbed}} = w_{\text{global}} + \epsilon_w$ and $\delta_j^{\text{perturbed}} = \delta_j + \epsilon_{\delta,j}$, we perform a second forward-backward pass using the perturbed model to evaluate the perturbed loss $L^{\text{perturbed}} = L(w_{\text{global}}^{\text{perturbed}}, \delta^{\text{perturbed}})$. 

We then update our active parameters using the preconditioned gradients of the perturbed loss, scaled by the Entropy-Adaptive Learning Rate $\eta_t$:
\begin{equation}
w_{\text{global}} \leftarrow w_{\text{global}} - \eta_t \frac{\partial L^{\text{perturbed}}}{\partial w_{\text{global}}}
\end{equation}
\begin{equation}
\delta_j \leftarrow \delta_j - \eta_t \frac{1}{\tilde{F}_j + \epsilon_{\text{stab}}} \frac{\partial L^{\text{perturbed}}}{\partial \delta_j}
\end{equation}
where:
\begin{equation}
\eta_t = \frac{\eta_0}{1.0 + \gamma_{\text{ealr}} H_{\text{avg}}}
\end{equation}
In our SOTA configuration, we set the base learning rate $\eta_0 = 0.1$, the EALR scaling factor $\gamma_{\text{ealr}} = 5.0$, the preconditioning damping constant $\epsilon_{\text{stab}} = 0.001$, and the SAM neighborhood size $\rho = 0.05$.

\end{document}
%%%%%%%%%% END OF FILE %%%%%%%%%%%%%%%%%
